% GENERATED FILE. Edit canonical lessons, then run npm run generate:books.
% canonical-source-hash: fnv1a64:95d75066ee6a1359
% canonical-lessons: SA-C336-nitih, SA-C336-lobhah, SA-C336-makarah, SA-C336-mudhah, SA-C336-lobhayati

\chapter{Fables}
\label{ch:sa-fables}
\hlchaptermodality{336}

\begin{chapteropening}
\textbf{By the end of this chapter:} \emph{I can tell a fable and say what it teaches.}

Complete the last lesson of chapter 336: I can tell a fable and say what it teaches.
\end{chapteropening}

\section[\texorpdfstring{\sk{नीतिः} (nītiḥ)}{नीतिः (nītiḥ)}]{\sk{नीतिः} (nītiḥ) — wise conduct; a moral}

\label{lesson:SA-C336-nitih}



\begin{quote}
Before the new one: say the Sanskrit for bewilders, then the Sanskrit for attracts.
\end{quote}

\subsection*{You'll want to know: \sk{नीतिः}}
\textbf{\sk{नीतिः}} — \emph{nītiḥ} — \textquotedblleft{}wise conduct; a moral\textquotedblright{}.

\begin{quote}
\sk{पितामहः} \sk{बालाय} \sk{नीतिं} \sk{कथयति।}
\end{quote}

\begin{quote}
\emph{pitāmahaḥ bālāya nītiṃ kathayati.}
\end{quote}

\begin{quote}
Grandfather tells the boy about wise conduct.
\end{quote}

A feminine noun.

\begin{cousinweb}
\textbf{\sk{नीतिः}} (\emph{nītiḥ}) comes from the root √nī, \textquotedblleft{}to lead\textquotedblright{}, as in \emph{nayati}; MW gives \emph{nīti} as guidance, management, and right or wise conduct, prudence (MW).
\end{cousinweb}

\begin{grammarlens}[title={What you've built}]
One more word for telling what happened.
\end{grammarlens}

\subsection*{Your turn}
\begin{itemize}
\raggedright
  \item \emph{Say it:} \emph{nītiḥ}
  \item \emph{Say it:} \emph{nītiḥ}, once more
  \item \emph{Say it:} \emph{pitāmahaḥ bālāya nītiṃ kathayati.} — Grandfather tells the boy about wise conduct.
  \item \emph{Recall:} say the Sanskrit for bewilders, then the Sanskrit for attracts, then say \emph{nītiḥ} again
\end{itemize}

\begin{tcolorbox}[breakable,colback=teal!4,colframe=teal!35!black,title={Before you move on}]
What does \sk{नीतिः} mean? (\textquotedblleft{}Wise conduct; a moral\textquotedblright{}.) Say it once more.
\end{tcolorbox}
\section[\texorpdfstring{\sk{लोभः} (lobhaḥ)}{लोभः (lobhaḥ)}]{\sk{लोभः} (lobhaḥ) — greed}

\label{lesson:SA-C336-lobhah}



\begin{quote}
Before the new one: say the Sanskrit for attracts, then the Sanskrit for wise conduct.
\end{quote}

\subsection*{You'll want to know: \sk{लोभः}}
\textbf{\sk{लोभः}} — \emph{lobhaḥ} — \textquotedblleft{}greed\textquotedblright{}.

\begin{quote}
\sk{लोभः} \sk{दुःखस्य} \sk{कारणम्} \sk{अस्ति।}
\end{quote}

\begin{quote}
\emph{lobhaḥ duḥkhasya kāraṇam asti.}
\end{quote}

\begin{quote}
Greed is a cause of sorrow.
\end{quote}

A masculine noun.

\begin{cousinweb}
\textbf{\sk{लोभः}} (\emph{lobhaḥ}) comes from the root √lubh, \textquotedblleft{}to be perplexed, desire greatly\textquotedblright{}; MW gives \emph{lobha} as confusion, eager desire, and covetousness, greed (MW).
\end{cousinweb}

\begin{grammarlens}[title={What you've built}]
One more word for telling what happened.
\end{grammarlens}

\subsection*{Your turn}
\begin{itemize}
\raggedright
  \item \emph{Say it:} \emph{lobhaḥ}
  \item \emph{Say it:} \emph{lobhaḥ}, once more
  \item \emph{Say it:} \emph{lobhaḥ duḥkhasya kāraṇam asti.} — Greed is a cause of sorrow.
  \item \emph{Recall:} say the Sanskrit for attracts, then the Sanskrit for wise conduct, then say \emph{lobhaḥ} again
\end{itemize}

\begin{tcolorbox}[breakable,colback=teal!4,colframe=teal!35!black,title={Before you move on}]
What does \sk{लोभः} mean? (\textquotedblleft{}Greed\textquotedblright{}.) Say it once more.
\end{tcolorbox}
\section[\texorpdfstring{\sk{मकरः} (makaraḥ)}{मकरः (makaraḥ)}]{\sk{मकरः} (makaraḥ) — a sea monster; a crocodile}

\label{lesson:SA-C336-makarah}



\begin{quote}
Before the new one: say the Sanskrit for wise conduct, then the Sanskrit for greed.
\end{quote}

\subsection*{You'll want to know: \sk{मकरः}}
\textbf{\sk{मकरः}} — \emph{makaraḥ} — \textquotedblleft{}a sea monster; a crocodile\textquotedblright{}.

\begin{quote}
\sk{नद्याम्} \sk{एकः} \sk{मकरः} \sk{वसति।}
\end{quote}

\begin{quote}
\emph{nadyām ekaḥ makaraḥ vasati.}
\end{quote}

\begin{quote}
A crocodile lives in the river.
\end{quote}

A masculine noun.

\begin{cousinweb}
\textbf{\sk{मकरः}} (\emph{makaraḥ}) MW gives \emph{makara} as a kind of sea monster, sometimes confused with the crocodile, shark or dolphin, without a derivation, so its origin is uncertain (MW).
\end{cousinweb}

\begin{grammarlens}[title={What you've built}]
One more word for telling what happened.
\end{grammarlens}

\subsection*{Your turn}
\begin{itemize}
\raggedright
  \item \emph{Say it:} \emph{makaraḥ}
  \item \emph{Say it:} \emph{makaraḥ}, once more
  \item \emph{Say it:} \emph{nadyām ekaḥ makaraḥ vasati.} — A crocodile lives in the river.
  \item \emph{Recall:} say the Sanskrit for wise conduct, then the Sanskrit for greed, then say \emph{makaraḥ} again
\end{itemize}

\begin{tcolorbox}[breakable,colback=teal!4,colframe=teal!35!black,title={Before you move on}]
What does \sk{मकरः} mean? (\textquotedblleft{}A sea monster; a crocodile\textquotedblright{}.) Say it once more.
\end{tcolorbox}
\section[\texorpdfstring{\sk{मूढः} (mūḍhaḥ)}{मूढः (mūḍhaḥ)}]{\sk{मूढः} (mūḍhaḥ) — bewildered; foolish}

\label{lesson:SA-C336-mudhah}



\begin{quote}
Before the new one: say the Sanskrit for greed, then the Sanskrit for a sea monster.
\end{quote}

\subsection*{You'll want to know: \sk{मूढः}}
\textbf{\sk{मूढः}} — \emph{mūḍhaḥ} — \textquotedblleft{}bewildered; foolish\textquotedblright{}.

\begin{quote}
\sk{मूढः} \sk{मकरः} \sk{वानरं} \sk{विश्वसिति।}
\end{quote}

\begin{quote}
\emph{mūḍhaḥ makaraḥ vānaraṃ viśvasiti.}
\end{quote}

\begin{quote}
The foolish crocodile trusts the monkey.
\end{quote}

\begin{cousinweb}
\textbf{\sk{मूढः}} (\emph{mūḍhaḥ}) comes from the root √muh, \textquotedblleft{}to be confused, go astray\textquotedblright{}; MW gives \emph{mūḍha} as stupefied, bewildered, perplexed, confused (MW).
\end{cousinweb}

\begin{grammarlens}[title={What you've built}]
One more word for telling what happened.
\end{grammarlens}

\subsection*{Your turn}
\begin{itemize}
\raggedright
  \item \emph{Say it:} \emph{mūḍhaḥ}
  \item \emph{Say it:} \emph{mūḍhaḥ}, once more
  \item \emph{Say it:} \emph{mūḍhaḥ makaraḥ vānaraṃ viśvasiti.} — The foolish crocodile trusts the monkey.
  \item \emph{Recall:} say the Sanskrit for greed, then the Sanskrit for a sea monster, then say \emph{mūḍhaḥ} again
\end{itemize}

\begin{tcolorbox}[breakable,colback=teal!4,colframe=teal!35!black,title={Before you move on}]
What does \sk{मूढः} mean? (\textquotedblleft{}Bewildered; foolish\textquotedblright{}.) Say it once more.
\end{tcolorbox}
\section[\texorpdfstring{\sk{लोभयति} (lobhayati)}{लोभयति (lobhayati)}]{\sk{लोभयति} (lobhayati) — tempts; lures}

\label{lesson:SA-C336-lobhayati}



\begin{quote}
Before the new one: say the Sanskrit for a sea monster, then the Sanskrit for bewildered.
\end{quote}

\subsection*{You'll want to know: \sk{लोभयति}}
\textbf{\sk{लोभयति}} — \emph{lobhayati} — \textquotedblleft{}tempts; lures\textquotedblright{}.

\begin{quote}
\sk{शृगालः} \sk{फलेन} \sk{काकं} \sk{लोभयति।}
\end{quote}

\begin{quote}
\emph{śṛgālaḥ phalena kākaṃ lobhayati.}
\end{quote}

\begin{quote}
The jackal tempts the crow with a fruit.
\end{quote}

Past (he …): \sk{सः} \sk{काकं} \sk{लोभितवान्}, \emph{saḥ kākaṃ lobhitavān}.

\begin{cousinweb}
\textbf{\sk{लोभयति}} (\emph{lobhayati}) comes from the root √lubh, \textquotedblleft{}to desire greatly\textquotedblright{}, as in \emph{lobha}, in its causative form: to make desire, tempt, lure (MW).
\end{cousinweb}

\begin{grammarlens}[title={What you've built}]
One more word for telling what happened.
\end{grammarlens}

\subsection*{Your turn}
\begin{itemize}
\raggedright
  \item \emph{Say it:} \emph{lobhayati}
  \item \emph{Say it:} \emph{lobhayati}, once more
  \item \emph{Say it:} \emph{śṛgālaḥ phalena kākaṃ lobhayati.} — The jackal tempts the crow with a fruit.
  \item \emph{Recall:} say the Sanskrit for a sea monster, then the Sanskrit for bewildered, then say \emph{lobhayati} again
\end{itemize}

\begin{tcolorbox}[breakable,colback=teal!4,colframe=teal!35!black,title={Before you move on}]
What does \sk{लोभयति} mean? (\textquotedblleft{}Tempts; lures\textquotedblright{}.) Say it once more.
\end{tcolorbox}
